\section{The learning rule, the reduced flow, and the proofs}
\label{app:derivation}

\subsection{Derivation of the update}

The receiver's belief before an interaction is the product of Gaussians
\eqref{eq:belief}. It observes $\sigma_e$ and treats it as a report of an unknown
true label $\sigma_T$ through a channel that flips with probability
$\varepsilon(z)$. Marginalising over labels and over the issue as the emitter may
have parsed it,
%
\begin{equation}
  L(\sigma_e\mid \hat{x}, u) =
  \sum_{\sigma_T}\int P(\sigma_e\mid \sigma_T, y, \hat{x}, u)\,
  P(\sigma_T\mid y, \hat{x}, u)\,P(y\mid \hat{x}, u)\,\mathrm{d}y ,
\end{equation}
%
with $P(\sigma_e\mid\sigma_T,z)=(1-\varepsilon(z))\delta_{\sigma_e,\sigma_T} +
\varepsilon(z)\delta_{\sigma_e,-\sigma_T}$ and
$P(\sigma_T\mid y,w)=\Theta(\sigma_T\,y\cdot w)$. Two modelling choices keep the
algebra Gaussian: taking $\varepsilon(z)=\Phi(z)$ with $z$ Gaussian of mean
$\mu_{e|r}$ and variance $V_{e|r}$ makes the flip probability a bounded function
of a Gaussian variable while preserving conjugacy of the affective sector, and
taking $P(y\mid\hat{x},u)$ Gaussian with variance $\|w\|^{-2}$ lets an
inexperienced agent be correspondingly unsure it has parsed the issue as the
emitter did. Averaging $L$ over the current belief collapses the Gaussian
integrals onto the two scaled fields and gives \eqref{eq:evidence}: the evidence
depends on the receiver's state only through $(\hw,\hmu)$. That is the origin of
the sector symmetry in Proposition~\ref{prop:sym}(iii).

The exact posterior is not Gaussian. Replacing it by the maximum-entropy
Gaussian relative to the prior, subject to matching the first two moments, gives
updates in terms of derivatives of $\log Z$, which in the scaled fields are
\eqref{eq:update-ideo}--\eqref{eq:update-aff} with the modulation functions
\eqref{eq:modulation}.

\subsection{The exact field increments, and why the reduced flow is legitimate}
\label{app:reduced}

The updates move $\hat{w}_r$ and $\mu_{e|r}$, while $\hw$ and $\hmu$ are scaled
versions whose normalisations also move. Carrying the algebra through for a
channel held at fixed $(\hat{x},\sigma_e)$, and writing
$a=\hat{x}\!\cdot\!C_r\hat{x}/\gamma_C^{2}$ and $b=V_{e|r}/\gamma_V^{2}$, both in
$(0,1)$:
%
\begin{equation}
  \Delta\hw = a\Bigl[\Fw - \tfrac{a}{2}\,\hw\,\FC\Bigr],
  \qquad
  \Delta\hmu = b\Bigl[\Fmu - \tfrac{b}{2}\,\hmu\,\FV\Bigr],
  \label{eq:exact-increments}
\end{equation}
%
with all modulation functions evaluated at $(\Hf,\hmu)$. Three observations
justify stating Proposition~\ref{prop:sym} for the leading flow.

\emph{The correction terms share the symmetry.} $\FC$ is even and $\hw$ odd
under $u\mapsto-u$, so $\hw\FC$ is odd, exactly like $\Fw$. Nothing in
Proposition~\ref{prop:sym}(iii) depends on dropping them.

\emph{They are cubic at the fixed point.} On the diagonal
$\Fw(s,s)=-(2/\pi)s+O(s^3)$ and $\FC(s,s)=O(s^2)$, so $\hw\FC=O(s^3)$. The
linearisation, the eigenvalues and the tangent to the separatrix are unaffected.

\emph{The isotropic case is the model's own initial condition.} With $C=I$,
$V=1$ and unit issues, $a=b=1/2$ exactly. Anisotropy $a\neq b$ tilts the
separatrix to slope $\sqrt{b/a}$ but preserves the hyperbolic saddle, so every
conclusion survives by structural stability. Part~(i) needs no assumption at
all: $D$ enters only through $\Hf$, in \eqref{eq:exact-increments} as much as in
the leading flow.

\subsection{Proof of Proposition~\ref{prop:sym}}

\emph{(i)} $D$ appears in $\Fw$ and $\Fmu$ only through $\Hf=\hw+D$, so
$X_D(\hw,\hmu)=X_0(\hw+D,\hmu)$ by inspection.

\emph{(ii)} $\Fw$ carries the factor $1-2\Phi(\hmu)$, which vanishes exactly at
$\hmu=0$, and $\phi/Z$ is finite and positive; $\Fmu$ carries $1-2\Phi(\Hf)$,
vanishing exactly at $\Hf=0$. The two nullclines meet only at $(-D,0)$.

\emph{(iii)} Take $D=0$. From $\Phi(-z)=1-\Phi(z)$,
$Z(-\hw,-\hmu)=(1-\Phi(\hw))+(1-\Phi(\hmu))-2(1-\Phi(\hw))(1-\Phi(\hmu))=Z(\hw,\hmu)$,
while $1-2\Phi(-\hmu)=-(1-2\Phi(\hmu))$ and $\phi$ is even. Hence
$\Fw(-u)=-\Fw(u)$ and likewise for $\Fmu$, so $X_0(-u)=-X_0(u)$. Sector exchange
is immediate from the symmetry of $Z$ in its two arguments, which gives
$\Fw(x,y)=\Fmu(y,x)$.

\emph{(iv)} On $\hmu=\Hf$ the two modulation functions coincide, since $Z$ is
symmetric and the prefactors exchange, so $\Fw=\Fmu$ there and the flow is
parallel to the line: it is invariant, not merely tangent. (On the antidiagonal
$\hmu=-\Hf$ one gets $\Fmu=-\Fw$, so that line is invariant too.) Linearising at
$(-D,0)$, the diagonal terms vanish by (ii) and
$\partial_{\hmu}\Fw|_0=\partial_{\Hf}\Fmu|_0=-2\phi(0)^2/Z(0,0)=-2/\pi$, using
$Z(0,0)=1/2$ and $\phi(0)^2=1/(2\pi)$. The Jacobian is
$\bigl(\begin{smallmatrix}0&-2/\pi\\-2/\pi&0\end{smallmatrix}\bigr)$, with
eigenvalues $\mp2/\pi$ and eigenvectors along the diagonal (stable) and
antidiagonal (unstable). \qed

\subsection{The gap distribution and $\Gresp$}
\label{app:gap}

At initialisation $\hat{w}\sim\mathcal{N}(0,I_K)$ and $\hat{x}$ is a unit vector,
so $\hat{w}_r\!\cdot\!\hat{x}\sim\mathcal{N}(0,1)$ and, with $C=I$ giving
$\gamma_C=\sqrt2$, $\hw\sim\mathcal{N}(0,\tfrac12)$. With $V=1$,
$\gamma_V=\sqrt2$ and $\mu\sim\mathcal{U}(-1,1)$, $\hmu$ is uniform on
$(-c,c)$ with $c=1/\sqrt2$. The gap $\psi=\hmu-\hw$ is therefore the convolution
of $\mathcal{U}(-c,c)$ with $\mathcal{N}(0,\sigma^2)$, $\sigma=1/\sqrt2$. Using
the antiderivative $\mathcal{F}(z)=z\Phi(z)+\phi(z)$ of $\Phi$,
%
\begin{equation}
  \Psi(t) \;=\; \frac{\sigma}{2c}\Bigl[\mathcal{F}\bigl(\tfrac{t+c}{\sigma}\bigr)
  - \mathcal{F}\bigl(\tfrac{t-c}{\sigma}\bigr)\Bigr],
  \qquad
  \Gresp(D)=2\Psi(D)-1 .
\end{equation}
%
$\Gresp$ is odd, strictly increasing and tends to $\pm1$; its slope at the origin
is $\Gresp'(0)=2p_\psi(0)=\bigl[\Phi(c/\sigma)-\Phi(-c/\sigma)\bigr]/c=0.9655$.
The accompanying code checks $\Psi$ against $2\times10^{6}$ Monte Carlo samples
to within $2\times10^{-3}$.

\subsection{The channel asymmetry of Corollary~\ref{cor:ordering}}

$\mathrm{sign}(\Fmu)=\mathrm{sign}(1-2\Phi(\Hf))=-\mathrm{sign}(\Hf)$, so the
field can invert the trust update whenever $|\hw|<|D|$ and the signs oppose. By
contrast $\mathrm{sign}(\Fw)=\mathrm{sign}(1-2\Phi(\hmu))$ does not involve $D$
at all, so the direction of every opinion update is set by trust alone. To first
order in $d$, expanding $\Fw(\hw+D,\cdot)=\Fw+D\,\FC+O(D^2)$ and averaging over
the initial measure at a class-symmetric state with $C_r=cI$, the
class-correlated drift of the opinion overlap is proportional to
$\mathbb{E}[\FC\,|w_e\!\cdot\!\hat{x}|]$, which evaluates to $-0.0247$: negative,
and some twenty-five times smaller in magnitude than the trust-channel
coefficient $-2/\pi$. Hence $\CcO$ is driven \emph{negative} at small $\fd|d|$ and
turns positive only when the trust field is strong enough to reorganise opinions
--- a level set of $\fd\Gresp(d)$, which is approximately hyperbolic at small $d$
and flattens as $\Gresp$ saturates.

\subsection{Why no critical field strength exists for $\CcT$}
\label{app:no-critical}

$\Fmu=(1-2\Phi(\Hf))\phi(\hmu)/Z$ carries a factor that vanishes identically on
$\Hf=0$. Hence $\Fmu\equiv0$ along that entire line and so, in particular,
$\partial_{\hmu}\Fmu|_{\Hf=0}=0$: linearising the affective sector about the
class-symmetric state gives a Jacobian with a zero eigenvalue in the
class-contrast direction. The contrast obeys
$\tfrac{\mathrm{d}}{\mathrm{d}\tau}(\mu_{\text{out}}-\mu_{\text{in}}) =
J_\mu(\mu_{\text{out}}-\mu_{\text{in}}) + b(\fd,d)$ with $J_\mu=0$, i.e. a driven
response with no restoring force, not an eigenvalue crossing. Together with the
class-relabelling symmetry --- which forces $\CcT$ to be odd in $d$ and exactly
zero at $d=0$ for every $N$ --- this establishes that the sign change at $d=0$
cannot be a bifurcation. A finite-size collapse performed around it would be a
collapse about a smooth odd function and would manufacture an exponent with no
referent.
